\documentclass[10.5pt,a4paper]{ctexart}

\usepackage[a4paper,margin=1.55cm]{geometry}
\usepackage{longtable,booktabs,array,tabularx}
\usepackage{xcolor}
\usepackage{enumitem}
\usepackage{hyperref}
\usepackage{graphicx}
\usepackage{titlesec}
\usepackage{multicol}
\usepackage{amssymb}
\usepackage{pifont}
\usepackage{fancyhdr}
\usepackage{lastpage}
\usepackage{float}
\usepackage{caption}

\hypersetup{
  colorlinks=true,
  linkcolor=blue!55!black,
  urlcolor=blue!55!black,
  citecolor=blue!55!black
}

\setlist[itemize]{leftmargin=1.45em,itemsep=0.17em,topsep=0.20em}
\setlist[enumerate]{leftmargin=1.55em,itemsep=0.17em,topsep=0.20em}
\renewcommand{\arraystretch}{1.18}
\setlength{\parskip}{0.30em}
\setlength{\parindent}{0pt}
\setlength{\LTpre}{0.3em}
\setlength{\LTpost}{0.3em}

\titleformat{\section}{\Large\bfseries\color{blue!55!black}}{\thesection}{0.6em}{}
\titleformat{\subsection}{\large\bfseries\color{black}}{\thesubsection}{0.6em}{}
\titleformat{\subsubsection}{\normalsize\bfseries\color{black}}{\thesubsubsection}{0.6em}{}

\pagestyle{fancy}
\fancyhf{}
\fancyhead[L]{拉萨--藏东南--藏北--羊湖--藏南大环线}
\fancyhead[R]{2026 实测修订通用版}
\fancyfoot[C]{第 \thepage\ 页 / 共 \pageref{LastPage} 页}
\setlength{\headheight}{14pt}

% ================== 信息框 ==================
\newcommand{\infobox}[3]{%
  \vspace{0.25em}
  \noindent\fcolorbox{#1!70!black}{#1!7}{%
    \parbox{0.965\linewidth}{\textbf{#2}\quad #3}}
  \vspace{0.25em}
}
\newcommand{\mainbox}[2]{\infobox{blue}{#1}{#2}}
\newcommand{\riskbox}[2]{\infobox{red}{#1}{#2}}
\newcommand{\methodbox}[2]{\infobox{green}{#1}{#2}}
\newcommand{\summarybox}[2]{\infobox{yellow}{#1}{#2}}
\newcommand{\tagmust}{\textcolor{green!40!black}{\textbf{强烈推荐}}}
\newcommand{\tagopt}{\textcolor{blue!70!black}{\textbf{可选}}}
\newcommand{\tagcut}{\textcolor{red!70!black}{\textbf{可砍}}}
\newcommand{\tagwarn}{\textcolor{red!70!black}{\textbf{高风险}}}
\newcommand{\ck}{\(\square\)}
\newcommand{\coord}[2]{\texttt{#1, #2}}
\newcommand{\grade}[1]{\textbf{#1/5}}

\title{\textbf{拉萨--藏东南--藏北--羊湖--藏南自驾大环线攻略}\\[0.3em]
\large 2026 年 8 月实测修订通用版：巴松措--林芝--鲁朗--易贡--尼屋--嘉黎--麦地卡--那曲--圣象天门--拉萨--羊卓雍措--卡若拉--普莫雍措--洛扎}
\author{基于一次完整自驾实测、OsmAnd 规划 GPX 与现场复盘整理}
\date{道路实测：2026-08-08 至 2026-08-14；政策信息核验至 2026-08-15}

\begin{document}
\maketitle

\mainbox{这份攻略和普通“景点清单”的区别}{这是一份\textbf{跑完以后倒推回来重写的执行型攻略}。景点评价、真实通行时间、烂路强度、补给节点、停车方式和砍点顺序，优先采用 2026 年 8 月实测；政策、预约和边境管理规则另外查验官方信息。GPX 为 OsmAnd 根据汽车/步行导航拼接生成的\textbf{规划轨迹，不是 GPS 实录轨迹}：共 67 个控制点、28628 个轨迹几何点，几何长度约 2608 km。}

\riskbox{先看这一句再决定是否抄路线}{\textbf{第一次上高原的人，不建议直接完整复制这条线。}路线会连续进入 4500--5200 m 高海拔，并包含多段长距离非铺装、坑洼县道和偏远山谷。若过去在 4000 m 左右有明显严重高反、驾驶非铺装经验不足、车辆为低底盘轿车，建议删掉“新措支线、易贡--尼屋穿越、东德错、藏南洛扎回程”中的至少一部分。}

\tableofcontents
\newpage

% ==========================================================
\section{先把路线看懂：这条环线到底在玩什么}

\subsection{景观主线不是“看一堆错”，而是连续跨越四种地貌}

这条路线最强的地方不是某一个景点，而是地貌连续变化非常明显：
\begin{enumerate}
  \item \textbf{拉萨河谷--尼洋河谷：}从高原城市进入藏东南，海拔总体下降，植被变多；米拉山口只是短暂拉高到 5000 m 级。
  \item \textbf{巴松措--鲁朗--易贡--尼屋：}森林、雪山、湖泊、深切峡谷和低海拔河谷，是全线最“湿润”的一段。
  \item \textbf{尼屋--嘉黎--麦地卡--那曲--纳木错：}森林峡谷逐步退出，地势抬升为开阔藏北高原，出现草甸、湿地、大湖和极强的尺度感。
  \item \textbf{拉萨--羊湖--卡若拉--普莫--洛扎：}从羊湖枝状湖汊进入冰川、雪山湖泊，再翻过蒙达拉进入洛扎亚高山河谷，最后由县道翻回羊湖东线。
\end{enumerate}

\summarybox{景观差异化}{巴松措看“雪山+森林+湖+草甸同框”；鲁朗看林海、田园和近距离雪山；易贡--尼屋看深切河谷；嘉黎--麦地卡看藏北高原湖泊湿地；圣象天门看纳木错的大尺度；羊湖看枝状湖汊和多个高度层次；卡若拉--恰央错--岗布沟看冰川近景；普莫--蒙达拉--洛扎看库拉岗日和藏南河谷。}

\subsection{推荐两种版本：7 天能跑，8--9 天更合理}

\begin{longtable}{p{2.3cm}p{5.8cm}p{5.9cm}p{2.6cm}}
\toprule
版本 & 适合谁 & 核心取舍 & 强度 \\
\midrule
\textbf{7 天高强度版} & 已适应高原、愿意每天长途驾驶、能接受临场砍点 & 基本复刻本次实测主线；巴松措只住一晚，新措可能与其他点冲突；圣象天门和最后藏南段时间都很紧 & 很高 \\
\textbf{8 天推荐版} & 大多数有高原经验的自驾游客 & 巴松措增加半天；林芝/鲁朗不赶夜路；圣象天门单独留足 4.5--5 h & 高但可控 \\
\textbf{9 天舒适版} & 想拍照、徒步、慢慢看雪山湖泊 & 巴松措 1.5--2 天；新措单独半天以上；圣象天门完整一天；羊湖/藏南不和赶飞机绑定 & 最推荐 \\
\bottomrule
\end{longtable}

\subsubsection{8--9 天推荐版建议拆法}
\begin{enumerate}
  \item D1 拉萨补给--米拉--太昭--巴松措/结巴村。
  \item D2 巴松措深玩：湖心岛、布多山栈道、遗忘码头、错高村北岸；若车辆和状态允许，新措支线单独占大半天。
  \item D3 巴松措--卡定沟--林芝补给--色季拉--鲁朗，住扎西岗村。
  \item D4 扎西岗/贡措--通麦--易贡--尼屋。
  \item D5 尼屋--独俊--错嘎错--嘉乃玉措--麦地卡/东德错--那曲。
  \item D6 那曲--蓬措--圣象天门--纳木湖镇--当雄。
  \item D7 羊八井--拉萨补给/西藏博物馆--羊湖 1--5 号观景台--羊湖西岸住宿。
  \item D8 卡若拉--恰央错--岗布沟--普莫雍措--蒙达拉--洛扎住宿（9 天版推荐在这里停）。
  \item D9 洛扎--Mai La 县道垭口--羊湖东线--贡嘎/拉萨。
\end{enumerate}

% ==========================================================
\section{出发前准备：车辆、高反、补给、证件}

\subsection{车辆门槛：完整跑线请把“能开”与“适合开”分开}

\begin{longtable}{p{4.5cm}p{2.3cm}p{3.8cm}p{5.5cm}}
\toprule
路线范围 & 最低建议 & 更稳妥 & 说明 \\
\midrule
拉萨--林芝--鲁朗--纳木错南北主线--羊湖主路 & 普通 SUV & 中型 SUV/四驱 & 大部分为铺装路，主要问题是高海拔、拥堵、坑洼和长距离 \\
加入巴松措普通景点、卡定沟、林芝城区 & 普通 SUV & 四驱 SUV & 无明显硬核越野需求 \\
\textbf{加入新措支线} & 高离地间隙 SUV，谨慎 & \textbf{四驱、厚胎壁、高离地间隙} & 实测/现场判断为极差颠簸路，轿车不建议 \\
\textbf{易贡--八盖--尼屋完整穿越} & 不建议普通城市 SUV 贸然尝试 & \textbf{四驱硬派越野或底盘防护充分的四驱 SUV} & 八盖--尼屋约 48 km 为本线最难路段，持续烂路并有爬升 \\
麦地卡--东德错 & 高离地间隙 SUV & 四驱 SUV & 最后约 10 km 非铺装，雨后风险显著上升 \\
圣象天门自驾 & 高离地间隙 SUV & 四驱 SUV & 官方也明确 S206 通往景区为一级土路、通行条件欠佳 \\
洛扎--Mai La--羊湖东线 & SUV & 四驱 SUV & 多坑路面，夜间不建议第一次跑 \\
\bottomrule
\end{longtable}

\methodbox{本次车辆实测}{使用长安深蓝 G318 增程 SUV。全程总计轻微托底约 4--5 次，底盘有钢板保护，因此未发生故障；轮胎幸运地未爆胎。5000 m 以上动力略有下降，但仍能完成爬升。这个结果说明“有底盘防护的四驱高底盘 SUV 可以跑”，\textbf{不等于低底盘城市 SUV 或轿车也应该照抄。}}

\subsection{路况 5 级定义}
\begin{longtable}{p{1.5cm}p{4.3cm}p{9.4cm}}
\toprule
等级 & 定义 & 执行含义 \\
\midrule
1/5 & 良好铺装 & 正常国省道，主要注意限速、落石、牲畜与天气 \\
2/5 & 铺装但窄/堵/有坑 & 普通 SUV 可走，预留时间，夜间慎行 \\
3/5 & 坑洼或短段非铺装 & 需要主动选线避坑，低底盘风险明显增加 \\
4/5 & 长距离烂路/强颠簸 & 四驱、高离地间隙、厚胎壁强烈推荐，雨天风险大 \\
4.5/5 & 极差支线 & 行驶时间远超导航，连续大坑/石块/车辙，轿车不建议 \\
5/5 & 本线最高难度 & 偏远、长距离持续烂路并伴随爬升，车辆、驾驶技术和天气三项缺一不可 \\
\bottomrule
\end{longtable}

\subsection{全线关键路况总表（2026 年 8 月实测）}
\begin{longtable}{p{4.0cm}p{1.5cm}p{2.5cm}p{7.2cm}}
\toprule
路段 & 评级 & 实测时间/尺度 & 关键提醒 \\
\midrule
拉萨--米拉--工布江达 & 1--2/5 & 正常铺装主路 & 第一日高反风险大于路况风险；米拉山口只短停 \\
结巴村--新措 & \textbf{4.5/5} & \textbf{单程约 2.5 h} & 全程强颠簸，往返驾驶约 5 h；整项至少预留 8 h；四驱高底盘强烈推荐 \\
林芝--色季拉--鲁朗 G318 & 2/5 & 导航时间易低估 & 道路狭窄、车流和大车多，会车慢；终年容易拥堵 \\
鲁朗--通麦 & 1.5/5 & 长下坡 & 约从 3300 m 降至 1900--2100 m，注意刹车热衰减、落石和湿滑 \\
易贡--八盖乡 & 2--2.5/5 & 多数已铺装硬化 & 仍有局部坑洼，注意躲坑 \\
\textbf{八盖乡--尼屋} & \textbf{5/5} & \textbf{约 48 km} & 全线最难。持续烂路+爬升，雨天/夜间不建议；车辆需底盘防护、四驱更稳 \\
尼屋--嘉黎 & 2/5 & 山区铺装为主 & 景观从森林峡谷向藏北高原切换，注意高海拔爬升 \\
麦地卡乡--东德错末段 & 3/5 & 末约 10 km 非铺装 & 从麦地卡入口到东德错约 1 h 15 min；比新措稍好，雨后另说 \\
那曲--圣象天门 & 1--4/5 & 最后约 50 km 烂路，约 1 h 20 min & 景区往返本身很耗时；进景区到出来建议总留 4.5--5 h \\
拉萨--羊湖 1--5 号 & 1--2/5 & 多观景台连续停车 & 六个点和路上拍照合计建议 2--2.5 h \\
卡若拉--普莫--洛扎主线 & 1--2.5/5 & 海拔高 & 景点徒步比驾驶更耗体力 \\
洛扎--Mai La--羊湖东线 & 3/5 & 洛扎回贡嘎约 5.5 h & 县道垭口和羊湖东线均有大量坑，避免夜间赶路 \\
\bottomrule
\end{longtable}

\riskbox{雨天特别规则}{藏东南（巴松措、林芝、鲁朗、易贡、尼屋）夏季降水明显多于藏北，本次巴松措夜间还遇到冰雹。\textbf{新措、八盖--尼屋、东德错三类非铺装/烂路支线，雨天必须重新评估，不要把晴天通行经验机械复制到暴雨后。}}

\subsection{高海拔：这条线有两个明显风险峰}

本线的高反风险不是平均分布的：
\begin{itemize}
  \item \textbf{第一风险峰：}抵达拉萨后的前 24 h + 米拉山口约 5000 m。身体尚未适应，所以第一天不建议再加思金拉错这类 4500 m 左右高海拔停留点。
  \item \textbf{第二风险峰：}嘉黎--麦地卡--那曲--圣象天门连续 4500 m 以上。到那曲后会有一天多持续高海拔暴露，必须降低活动强度。
  \item \textbf{第三类局部高点：}那根拉、卡若拉、恰央错、岗布沟等 5000 m 级景点。这里的问题不是短暂停车，而是\textbf{还要爬升/徒步}。
\end{itemize}

\riskbox{医疗底线}{高原急性病不能靠“硬扛”。Wilderness Medical Society 2024 指南强调：严重急性高山病、疑似高原脑水肿或高原肺水肿时，\textbf{下降海拔是核心处理}；不要继续上升。若出现意识改变、走路不稳、静息气短、明显进行性呼吸困难等危险表现，应立即停止行程并寻求医疗/救援。药物预防与治疗应遵医嘱，本攻略不提供个体化用药剂量。}

\subsubsection{实用物资经验}
\begin{itemize}
  \item 三人 7 天高原自驾，本次经验量级为\textbf{便携氧气约 8--10 瓶}；它只是临时辅助，不替代下降海拔和就医。
  \item 拉萨第一天一次补齐：氧气、饮用水、功能饮料、路餐、常用药；林芝再补一次，之后优质大型补给点明显变少。
  \item 3500 m 以上酒店建议优先选\textbf{有加湿器、供氧条件}的房间；4000 m 以上更推荐供氧房。高海拔极干，润唇膏、保湿霜和持续补水很重要。
  \item 高原上体力消耗和饥饿感会更明显。不要为了赶景点把午饭完全砍掉；至少准备能量棒、面包、肉干、巧克力和电解质饮料。
\end{itemize}

\subsection{证件与边境管理}
\begin{itemize}
  \item 身份证、驾驶证、车辆租赁/保险资料必须随身。
  \item 进入边境管理区前，\textbf{建议按规定办理边境管理区通行证}。2026 年 4 月起，16 周岁以上中国内地居民可通过“移民局12367”APP/微信/支付宝小程序申请电子边境管理区通行证，建议提前截图并保存离线副本。
  \item 本次实测在浪卡子以南区域遇到检查，现场主要核验身份证，出区又检查一次；\textbf{这只是一次实际检查口径，不代表通行证可以不办。}
  \item 西藏自治区边境管理条例明确：户籍不在边境管理区的中国公民进入边境管理区，应持居民身份证和边境管理区通行证。
\end{itemize}

\subsection{导航：高德负责主路，OsmAnd 负责“别人地图里没有的那条路”}
\begin{itemize}
  \item 高德/百度：城区、主干道、加油/充电、普通景区。
  \item OsmAnd：离线底图、县乡道、东德错支线、羊八井野温泉入口、Mai La 县道垭口、整条 GPX 轨迹。
  \item \textbf{不要把导航 ETA 当真实通行时间。}新措、八盖--尼屋、圣象天门末段就是典型反例。
  \item 偏远段遇到施工、塌方、临时管制，以现场交警/施工人员/酒店老板/当地司机的口径优先。
\end{itemize}

% ==========================================================
\section{点位导航总表：GPX 控制点、观景台与停车方式}

\mainbox{坐标怎么用}{以下坐标来自本次 OsmAnd 规划 GPX（WGS84），适合做\textbf{路线定位和岔路辅助}；它不是实录 GPS，也不是官方景区坐标数据库。常规景区优先直接搜索官方名称；只有野点、县道垭口、无标准中文名的位置才建议把坐标当备用。}

\begin{longtable}{p{3.4cm}p{3.6cm}p{4.0cm}p{5.1cm}}
\toprule
点位 & GPX/导航 & 停车 & 实用说明 \\
\midrule
贡嘎机场 & \coord{29.29597}{90.89638} & 机场/租车区正规停车 & 环线起止参考点 \\
拉萨城区锚点 & \coord{29.65254}{91.12209} & 不建议拿此坐标当停车位 & 市中心共享电动车存在大片禁停区，布宫/大昭寺/宇拓路尤其明显 \\
米拉山口 & \coord{29.82155}{92.36628} & 路侧正规停车区/观景台 & 第一日只短停，减少 5000 m 级暴露 \\
太昭古城 & \coord{30.00118}{93.07377} & 景区/古城附近停车 & 看尼洋河和古城氛围，价值主要在河谷而非古建本身 \\
巴松措景区入口一带 & \coord{30.00950}{93.91579} & 游客中心正规停车 & 进入湖区前后的补给/购票锚点 \\
巴松措湖区锚点 & \coord{30.00997}{93.97196}、\coord{30.04526}{94.02894} & 各观景点均有停车区域 & 结合达切拉、湖心岛、布多山、遗忘码头、错高村使用 \\
新措支线起/徒步端 & \coord{30.00333}{94.21345} -- \coord{30.02483}{94.23280} & 以现场停车点为准 & GPX 中为汽车+步行组合节点；支线极烂，务必单独留时间 \\
秀巴古堡 & \coord{29.84082}{93.71286} & 景区停车场 & 本次评价偏低，可直接砍 \\
卡定沟 & 直接导航“卡定沟景区” & 景区停车场 & 基本不绕路，45--60 min 快游足够 \\
比日神山观景区域 & GPX：\coord{29.64570}{94.38944} & 导航“比日神山观景台/生态景区” & 一定要导到观景位置；俯瞰感有，但看林芝市区/尼洋河并不清楚 \\
林芝补给锚点 & \coord{29.64673}{94.36251} & 商圈/公园周边正规停车 & 第二次关键补给；建议把晚饭也在林芝解决 \\
色季拉山口 & \coord{29.62015}{94.66621} & 山口观景停车区 & 南迦巴瓦受云雾影响极大，不要为等山耗死后续时间 \\
鲁朗林海观景台 & \coord{29.61566}{94.69763} & 观景台停车区 & 看森林和河谷；若天色晚可短停 \\
扎西岗/贡措一带 & GPX 约 \coord{29.75131}{94.73396} -- \coord{29.76046}{94.74144} & 村内/景点停车位 & 天气好时加拉白垒和一众 6000 m 雪山非常值得 \\
通麦 & \coord{30.10262}{95.08454} & 路侧/镇区正规停车 & 通麦特大桥约 \coord{30.09783}{95.06532}；无需久留 \\
易贡方向锚点 & GPX \coord{30.20182}{94.88998} & 湖岸/茶场常见停车位置 & 易贡错重点看雪山；随后进入狭长峡谷 \\
尼屋乡 & \coord{30.48581}{93.98430} & 乡镇停车 & 穿越终点；建议天黑前到 \\
嘉黎县 & \coord{30.64320}{93.24542} & 县城正规停车 & 午餐、燃油/补给、休息 \\
嘉乃玉措方向 & GPX \coord{30.62064}{93.18025} & 湖边观景停车点 & 与巴松措气质不同，开阔高原湖景，值得停 \\
麦地卡/东德错方向 & GPX \coord{30.98565}{92.94729} & 湿地道路/湖边安全停车区 & 保护区内不要驶离既有道路压草场；末段非铺装 \\
那曲 & \coord{31.46783}{92.03887} & 酒店/城区 & 4500 m 级住宿，高反管理重点 \\
圣象天门北线中途锚点 & \coord{31.08739}{90.79367} & 以现有停车带为准 & 蓬措支线可选择性停，别因小湖挤占圣象天门主时间 \\
圣象天门景区核心 & GPX \coord{30.84837}{90.61133}，内部多个点 \coord{30.81839}{90.61255}、\coord{30.83740}{90.49320}、\coord{30.82135}{90.50838} & 景区指定停车点 & 入口换票后按景区路线开，看到合适观景点就停；严禁驶入未开放草场 \\
纳木湖镇 & \coord{30.73658}{91.08743} & 镇区 & 可补电、补给 \\
那根拉/当雄方向 & GPX \coord{30.68113}{91.10022} -- \coord{30.48056}{91.09999} & 山口停车区/县城 & 那根拉本身一般，可短停或砍 \\
羊八井野温泉入口区域 & GPX \coord{30.06092}{90.49138} & 小路附近选择安全位置 & 蓝色天国东约 100 m 的小路进入约 3 km；只看不下水 \\
汤旺温泉/服务区 & \coord{29.98096}{90.35029} & 服务区式停车 & 人工温泉+充电，适合补能；本次到达但未泡 \\
拉萨柳梧/万达补给 & GPX \coord{29.64052}{91.10208} & 万达地下停车 & 吃饭+永辉补给一站解决 \\
罗布林卡/藏博区域 & GPX \coord{29.65393}{91.09139} & 罗布林卡停车场/藏博停车 & 两处相邻，尽量一次停车步行串联 \\
雅江河谷观景台 & GPX \coord{29.22933}{90.62528} & 观景台停车 & 羊湖前的河谷过渡点 \\
羊湖 1--5 号观景带 & GPX 约 \coord{29.19415}{90.61792}、\coord{29.18778}{90.61720}、\coord{29.18760}{90.58494}、\coord{29.15160}{90.49410}、\coord{29.12514}{90.43855} & 各正规观景台 & 1 高位俯瞰；2 河谷俯瞰；3--4 近景，4 可到湖边；5 看远处雪山 \\
卡若拉徒步点 & GPX \coord{28.89919}{90.16410} -- \coord{28.90052}{90.16952} & 景区停车 & 栈道爬升近看冰舌；高海拔慢走 \\
恰央错徒步点 & GPX \coord{28.89850}{90.22517} -- \coord{28.89228}{90.22521} & 景区/路侧停车 & 可环湖体会雪山脚下开阔山谷 \\
岗布沟徒步点 & GPX \coord{28.92763}{90.24564} -- \coord{28.92886}{90.23319} & 景区停车 & 2026-08 实测栈道受塌方损坏，康布措无法正常抵达，价值打折 \\
普莫雍措 & GPX \coord{28.62320}{90.44974} & 观景点停车 & 看大湖和远处边境雪山；不要只在路边拍一张就走 \\
蒙达拉/库拉岗日方向 & GPX \coord{28.41904}{90.55786} -- \coord{28.41279}{90.57296} & 垭口/观景停车点 & 翻垭口后库拉岗日近景明显增强 \\
洛扎县 & \coord{28.38931}{90.85878} & 县城 & 9 天版推荐在这里住一晚 \\
Mai La 县道垭口 & GPX \coord{28.52298}{90.88540} & 无标准景区停车场，安全宽处短停 & \textbf{OsmAnd 名称“Mai La”，非正式中文景区名}；路面坑多 \\
羊湖东线回程锚点 & \coord{28.94651}{91.05613} & 观景停车带 & 坑洼较多，预留时间 \\
贡嘎回程锚点 & \coord{29.27682}{91.09860} -- 机场 & 正规停车 & 进入机场前提前安排还车/加油/充电 \\
\bottomrule
\end{longtable}

% ==========================================================
\section{逐日详解：推荐版怎么玩、看什么、哪里该砍}

\subsection{D1：拉萨补给--米拉山口--太昭古城--巴松措}

\subsubsection{核心任务}
第一天真正重要的不是“多打卡”，而是三件事：\textbf{完成高原落地适应、把补给一次买齐、天黑前进入巴松措住宿区。}

\begin{itemize}
  \item 拉萨城区建议把午饭、氧气、水、功能饮料、路餐、常用药、车上零食一次解决。
  \item 若只是看布达拉宫外观，可在药王山/白塔区域短停步行；本次未自驾到此，因此不提供未经实测的停车位硬推荐。
  \item \textbf{第一天不建议加入思金拉错。}它海拔约 4500 m，和“第一天快速适应、米拉短停后下撤”的目标冲突。
  \item 米拉山口约 5000 m，只停 5--10 min 拍照即可，不要在第一天进行大强度活动。
\end{itemize}

\subsubsection{太昭古城：为什么保留但只给 20--40 min}
本次认为太昭古城真正有价值的是\textbf{尼洋河河谷景观和小镇氛围}。后面林芝城区缺少一个同样顺路又舒服的尼洋河高质量观景台，所以可以在这里先完成尼洋河观景。古城本身无需投入太长时间。

\subsubsection{巴松措：全线最值得加时长的地方之一}
\tagmust。若有 8--9 天，推荐至少给 1.5 天；摄影/徒步型游客可给 2 天。

推荐点位：
\begin{itemize}
  \item \textbf{达切拉观光台：}进入湖区后的高位总览。
  \item \textbf{湖心岛：}经典景观，停车方便，但游客最多。
  \item \textbf{布多山观景栈道：}获得更高视角。
  \item \textbf{遗忘码头：}适合湖岸氛围和早晚光线。
  \item \textbf{错高村及其前方 2--3 km 观景台：}北岸视角好，游客相对分散。
  \item \textbf{去新措路上：}森林、溪流和山谷本身就有价值，不要只把它当“去新措的烂路”。
\end{itemize}

\methodbox{住宿}{住结巴村最方便。临湖房车/露营类酒店的优点是早晨直接看湖、日出和云海；本次实测房间较小、卫生间和整体设施偏简陋。若重视睡眠和高原恢复，宁可牺牲一点“开门见湖”换更好的房间、供氧和加湿器。}

\subsection{D2：巴松措/新措--卡定沟--林芝--色季拉--鲁朗}

\subsubsection{新措：景点很值，但交通比想象中狠得多}
\tagmust（有合适车辆和完整时间时）。本次因为时间冲突最终未进入，因此\textbf{景观推荐来自前期规划与当地路线信息，路况和通行时间按现场实际反馈修订，并非本次亲自开完整支线。}

\begin{itemize}
  \item 结巴村到新措单程约 \textbf{2.5 h}，不是原先估计的 50 min。
  \item 路况评级 \grade{4.5}：连续颠簸、坑洼和非铺装感很强。
  \item 往返纯驾驶约 5 h；加徒步、拍照和休息，\textbf{至少给 8 h}。
  \item 四驱越野/高底盘四驱 SUV 强烈推荐；城市 SUV 慎行，轿车不建议。
  \item 如果前一晚下暴雨/冰雹，第二天早上先问酒店和当地司机，不要盲目进入。
\end{itemize}

\subsubsection{秀巴古堡、卡定沟、比日神山怎么取舍}
\begin{itemize}
  \item \textbf{秀巴古堡：}\tagcut。本次是全线主观评价最低的景点之一，需要门票，除非特别喜欢古碉楼/历史，否则可直接省时间。
  \item \textbf{卡定沟：}\tagopt 偏推荐。基本不绕路，核心是峡谷和天佛瀑布，45--60 min 足够，与巴松措/鲁朗有一定差异化。
  \item \textbf{比日神山：}\tagopt。一定导航到观景台。可以俯瞰山地景观，但由于距离和遮挡，林芝市区与尼洋河并没有想象中清楚；时间紧时可以跳过。
\end{itemize}

\subsubsection{林芝补给要当成正式任务}
林芝是第二个关键补给点：第一天拉萨遗漏的氧气、水、路餐、药、车辆用品，应在这里补齐。\textbf{建议晚餐也在林芝解决}，再翻色季拉去鲁朗；否则抵达鲁朗往往已经很晚，晚饭选择和睡眠时间都会被压缩。

\subsubsection{色季拉与南迦巴瓦：降低预期反而容易开心}
色季拉理论上可以看南迦巴瓦，但只要有云雾就会挡，而且距离较远。把它当\textbf{天气好就中奖的山口}，不要为了等山拖到夜路。

\subsection{D3：鲁朗深玩--通麦--易贡--尼屋穿越}

\subsubsection{鲁朗：全线第二个值得住下来慢玩的节点}
推荐住扎西岗村。天气好时，往北去贡措湖一带可以更近地看到\textbf{加拉白垒峰和一众 6000 m 级雪山}，体感比前一天在色季拉看南迦巴瓦更近、更有压迫感。

\begin{itemize}
  \item 扎西岗村：田园、木屋、草甸、林海和雪山组合。
  \item 贡措湖：湖不大，但优势是与近雪山和森林一起看。
  \item 若前一天抵达鲁朗太晚，可把“鲁朗林海收费观景台”降级，第二天早晨优先扎西岗/贡措。
\end{itemize}

\subsubsection{鲁朗--通麦：长下坡，不要只盯风景}
海拔从约 3300 m 持续下降到约 1900--2100 m，是西藏少见的湿润河谷森林路段。景观有点类似江浙沪深山河谷，但尺度更深切；因为河谷狭长，反而不容易看到远处雪山。

驾驶重点：
\begin{itemize}
  \item 连续长下坡注意刹车温度，增程/混动车合理利用能量回收和低挡控制车速。
  \item 林芝--鲁朗--通麦 G318 车多、路窄，大车会车耗时，导航 ETA 要加缓冲。
\end{itemize}

\subsubsection{易贡--尼屋：这条线值不值，取决于你愿不愿意吃 48 km 烂路}
景观上它非常值得：
\begin{enumerate}
  \item 易贡错段重点是雪山与湖谷。
  \item 离开易贡错后进入狭窄深切峡谷，主要看溪流、森林和山体。
  \item 临近尼屋后河谷突然重新开阔，视觉变化很漂亮。
\end{enumerate}

路况分段：
\begin{itemize}
  \item \textbf{易贡--八盖乡：}基本完成铺装硬化，局部坑洼，躲坑即可。
  \item \textbf{八盖乡--尼屋约 48 km：}\tagwarn，\grade{5}。持续烂路、颠簸并有爬升，本次为全线最难段。
\end{itemize}

\riskbox{进入条件}{只建议在\textbf{白天 + 天气稳定 + 车辆状态良好 + 驾驶员有非铺装经验}时进入。不要在下午很晚才从易贡出发“赌能赶到”。如果前方施工/塌方信息不明确，宁可绕行或取消。}

\subsection{D4：尼屋--独俊大峡谷--错嘎错--嘉乃玉措--麦地卡--东德错--那曲}

\subsubsection{尼屋到嘉黎：真正好看的是“地貌变了”}
从尼屋离开后，森林峡谷逐步退场；经过独俊大峡谷的过渡，海拔上升、山势变缓、视野变得越来越开阔，开始进入典型藏北高原。

\begin{itemize}
  \item \textbf{独俊大峡谷：}适合作为地貌转换的节点看，不必长停。
  \item \textbf{错嘎错：}就在嘉黎县城东部一带，顺路，可短停。
  \item \textbf{嘉乃玉措：}\tagmust。本次认为很漂亮，气质与巴松措/新措明显不同：森林背景减少，湖和高原空间感更突出。
\end{itemize}

\subsubsection{麦地卡湿地与东德错}
麦地卡不是传统“景区大门+观光车”玩法，重点是高原湿地的广阔、河流和大量湖沼。国家级自然保护区内分布众多湖泊，也是重要水禽和野生动物栖息地。

\begin{itemize}
  \item 跟随 OsmAnd/既有道路进入东德错方向，麦地卡乡入口到湖边约 1 h 15 min。
  \item 最后约 10 km 非铺装，\grade{3}，比新措好，但雨后情况可能完全不同。
  \item \textbf{不要离开既有车辙/道路压草场，不追动物、不驶入保护区未开放区域。}
  \item 时间不足时，宁可只看麦地卡湿地主体，不必强行进东德错。
\end{itemize}

\subsubsection{那曲：不是“睡一觉就走”的普通县城节点}
那曲海拔约 4500 m，本线从这里开始进入第二个高反风险峰，并且次日还会去圣象天门。到酒店后别再安排夜游、跑步和长距离步行，优先吃饭、补水、供氧和睡觉。

\subsection{D5：那曲--蓬措--圣象天门--纳木湖镇--那根拉--当雄}

\subsubsection{圣象天门为什么值得单独留一天}
\tagmust，且是全线湖泊景观的第一梯队。本次主观评价：\textbf{纳木错圣象天门是这条路线所有湖泊里最值得专程跑的一处。}

从那曲过去很远，末段还烂，但景观尺度、湖水颜色、念青唐古拉背景和靠近湖边的体验明显高一个量级。

\subsubsection{当前自驾规则（2026-08-15 核验）}
2026 年 6 月 15 日起，圣象天门取消原“自驾提前三日预约”和“每日 50 辆”限制。游客通过官方渠道购票并核验人员、车辆信息后，可按景区指定路线自驾入园。官方同时明确：S206 通往景区道路为\textbf{一级土路、通行条件欠佳}，并提供自愿购买的观光车方案。

\subsubsection{实际怎么玩}
\begin{itemize}
  \item 景区大门/游客接待处先换票、核验。
  \item 从进入景区到退出，\textbf{总计至少留 4.5--5 h}。
  \item 进核心区前沿途就有多个很好的纳木错视角，看到正规停车点和好光线就停，不必只冲“圣象”那个石门。
  \item 官方要求按指定线路行驶、指定区域停车；不要因为想拍“无人湖岸”就压草场或进入未开放区域。
\end{itemize}

\subsubsection{蓬措、那根拉、纳木湖镇}
\begin{itemize}
  \item \textbf{蓬措：}\tagopt。去圣象天门途中可以沿小路靠近湖边，水色很蓝、人少；但圣象天门时间不够时首先砍它。
  \item \textbf{纳木湖镇：}非常实用的补电/补给节点。
  \item \textbf{那根拉山口：}\tagcut。海拔很高，但景观本身一般；已经看完圣象天门时，没必要为了“到此一游”长停。
\end{itemize}

\subsection{D6：当雄--羊八井野温泉--汤旺温泉--拉萨--羊湖 1--5 号}

\subsubsection{羊八井：别只导航到商业温泉}
本次实测的玩法分成两段：
\begin{enumerate}
  \item \textbf{野温泉地热景观：}蓝色天国景区（当时未营业）东约 100 m 有小路，进入约 3 km，可看到自然热泉，水温体感很高（约 60℃量级）。\textbf{不要下水}，只在安全位置观察地热景观。
  \item \textbf{汤旺温泉/服务区：}继续向西约 20 km，有人工温泉、室内外泡池和充电站。本次到此补电但未泡。适合作为“补能+休息”的服务区式节点。
\end{enumerate}

\subsubsection{拉萨：万达补给 + 西藏博物馆优先级最高}
\begin{itemize}
  \item 柳梧新区万达适合停车、吃午饭、永辉补给，一次完成。
  \item \textbf{西藏博物馆：}\tagmust。本次评价非常高，馆体大、内容多，建议优先看 1 号、8 号展馆，并至少留 1.5--2.5 h。
  \item \textbf{罗布林卡：}\tagopt。本次主观感受一般，对藏式宫苑/宗教文化兴趣不强的人，只看园林和外部建筑可能觉得 60 元票价性价比普通。时间不足先砍罗布林卡，不砍藏博。
  \item 罗布林卡和西藏博物馆相邻，尽量一次停车步行串联。罗布林卡停车场有较大停车容量；藏博也有院内/周边停车安排，但暑期可能拥堵。
\end{itemize}

\subsubsection{羊湖 1--5 号观景台：不要只打卡 1 号}
沿拉萨往羊湖西岸，会连续获得不同高度的湖景：
\begin{itemize}
  \item \textbf{1 号：}山顶垭口大俯瞰，最经典的“羊湖全景”。
  \item \textbf{2 号：}进入河谷后的中高位俯瞰，湖岸关系更清楚。
  \item \textbf{3 号：}河谷近景，开始感受到水面尺度。
  \item \textbf{4 号：}可以下到湖边，拍照/打卡设施更多。
  \item \textbf{5 号：}湖景叠加远处雪山，天气通透时可看宁金抗沙峰等雪山方向。
\end{itemize}

\summarybox{时间}{把雅江河谷观景台也算上，从进入羊湖观景带到西岸住宿，\textbf{建议给 2--2.5 h}，不要按“导航 1 h 20 min”计算。}

\subsection{D7：羊湖--卡若拉--恰央错--岗布沟--普莫--蒙达拉--洛扎--Mai La--贡嘎}

\riskbox{先说结论}{这是最不建议和“当天晚上赶飞机/必须还车”绑定的一天。7 天版能跑，但压力很大；\textbf{9 天版强烈建议在洛扎住一晚}，第二天再走 Mai La 和羊湖东线。}

\subsubsection{卡若拉冰川 + 恰央错 + 岗布沟：至少 3 h，慢的人 4--5 h}
三个点都在约 5000 m 海拔进行步行/爬升，真正限制不是里程，而是高海拔体力。

\begin{itemize}
  \item \textbf{卡若拉冰川：}\tagmust。沿栈道上到高处，可以更近看冰舌，同时看远处雪山。
  \item \textbf{恰央错：}\tagmust。重点是雪山脚下小湖和开阔山谷，可环湖走一段，氛围与羊湖/纳木错都不同。
  \item \textbf{岗布沟：}\tagopt。理论上康布措直接在雪山脚下，非常漂亮；但 2026-08 实测因塌方导致栈道损坏，无法正常抵达湖泊，只能远看冰川，体验明显打折。出发前必须重新确认维修情况。
\end{itemize}

\subsubsection{普莫雍措--蒙达拉--库拉岗日：真正的藏南高潮}
\begin{itemize}
  \item 普莫雍措先看大尺度远景，远处为中不边境方向雪山，但部分会被近山遮挡。
  \item 翻 G219 蒙达拉山垭口后，库拉岗日雪山群明显变近，是这段最重要的雪山视角。
  \item 继续下降进入洛扎河谷，向东可见形态很有辨识度的夏拉拉岗、打拉日“双子山”方向。
  \item 洛扎河谷植被明显增加：谷底有树木和绿色带，而上方山体干燥裸露，有点像高原版“绿洲”，与前面所有天都不一样。
\end{itemize}

\subsubsection{推瓦寺和日托寺：为什么通用版不把它们设成必去}
原计划包含推瓦寺、日托寺，但本次实际因为时间不够全部取消。通用版建议：
\begin{itemize}
  \item 推瓦寺只在前面卡若拉三景、普莫都没超时的情况下再加；否则蒙达拉/库拉岗日优先级更高。
  \item 日托寺会带来明显绕行和时间压力；若最终目的地是贡嘎机场，不建议为了“补景点”把回程拖到夜间。
\end{itemize}

\subsubsection{洛扎--Mai La--羊湖东线--贡嘎：回程 5.5 h 不是开玩笑}
本次实际从洛扎县直接北上翻县道垭口（OsmAnd 标记为 \textbf{Mai La}），接入羊湖东线，再回贡嘎机场，\textbf{总计约 5.5 h}。

\begin{itemize}
  \item Mai La 县道：坑很多，视线和避坑比速度重要。
  \item 羊湖东线：同样存在大量坑洼，不能拿铺装国道平均速度估算。
  \item 不熟悉路况时不要夜跑；如果前面超时，宁可在洛扎/浪卡子方向住宿，第二天再出。
\end{itemize}

% ==========================================================
\section{停车、补给与住宿：最容易被“攻略省略”的执行细节}

\subsection{停车总原则}
\begin{itemize}
  \item 大多数自然景点都有明确观景台或停车区域，\textbf{优先停现成硬化区域}，不要为了机位把车压到草场、湿地和私人土地。
  \item 圣象天门必须按景区指定线路和停车管理执行。
  \item 羊湖 1--5 号均按现有观景台停车，不要在窄弯道突然刹停。
  \item 拉萨城区是全线唯一明显“停车比景点本身更麻烦”的地方：万达、罗布林卡/藏博分别解决，不建议频繁横穿老城找位置。
  \item 药王山/布宫本次未自驾停车，不提供假装实测的停车建议；如果朋友一定要开车去，建议出发当天在导航里搜正规停车场并预留拥堵时间。
\end{itemize}

\subsection{补给优先级}
\begin{longtable}{p{3.2cm}p{3.0cm}p{8.4cm}}
\toprule
地点 & 优先级 & 建议补什么 \\
\midrule
拉萨第一天 & \textbf{最高} & 氧气、水、功能饮料、路餐、药品、防晒/保湿、车内零食 \\
林芝 & \textbf{最高} & 把第一天遗漏全部补完；建议顺便吃晚饭 \\
鲁朗 & 中 & 餐饮、简单水和零食；不要指望像林芝一样全面 \\
嘉黎 & 中高 & 午餐、燃油/基本补给；进麦地卡前检查水和油 \\
那曲 & 高 & 高海拔住宿补给、氧气、晚餐、第二天路餐 \\
纳木湖镇 & 高 & 补电/补给，圣象天门后非常实用 \\
汤旺温泉服务区 & 中高 & 补电、休息，可泡温泉 \\
拉萨柳梧万达 & \textbf{最高} & 正餐+永辉超市，一次解决羊湖/藏南前的最后大补给 \\
浪卡子/洛扎 & 中高 & 藏南段最后补水、吃饭、油量检查 \\
\bottomrule
\end{longtable}

\subsection{增程/新能源车经验}
本次 G318 实际油电合计约 1104 元量级（具体价格和里程会变），增程车在这种路线的优势是：县城/服务区能充就充，偏远段用油兜底。

\begin{itemize}
  \item 不为了便宜电价绕路；每天优先保证油箱有安全余量。
  \item 巴松措景区大门前有中石化充电站，可在进入/离开时补电。
  \item 那曲、纳木湖镇、当雄、羊八井/汤旺、拉萨等节点补能条件明显更好。
  \item 高海拔、烂路、低温/雨天会使能耗高于平原；不要按日常市区电耗推续航。
\end{itemize}

\subsection{住宿选择：这条线不要在房费上过度省}
\begin{itemize}
  \item 藏东南低海拔（林芝、易贡一带）：重点是干燥/潮湿舒适度、热水、停车。
  \item 拉萨/鲁朗：如果已经适应，高品质普通房即可；但加湿器仍非常有用。
  \item 嘉黎、那曲、当雄、羊湖等 4000 m 左右或以上：优先供氧、加湿器、暖气/地暖和好床垫。
  \item 偏远地区\textbf{睡眠恢复本身就是安全配置}，宁可少去一个鸡肋景点，也不要为了省 100--200 元住得一夜睡不着。
\end{itemize}

\subsubsection{本次住宿的一个典型取舍：巴松措湖边房车}
优点：临湖、早晨方便看日出和云海，氛围强。缺点：房间小、卫生间和设施一般。通用建议是根据同行人耐受度选择：摄影党会觉得值，重视睡眠/卫生/高原恢复的人可能更适合结巴村里条件更好的酒店。

% ==========================================================
\section{景点评价与砍点优先级：时间不够时先砍谁}

\begin{longtable}{p{3.3cm}p{2.0cm}p{7.3cm}p{3.0cm}}
\toprule
点位 & 推荐度 & 主要价值 & 时间不够时 \\
\midrule
巴松措整体 & 5/5 & 全线最完整的雪山+森林+湖+草甸组合 & 不砍，反而建议加时 \\
新措 & 4.5/5 & 深山森林溪流与湖泊，交通本身也有景 & 车辆/天气不合适就整项取消 \\
太昭古城 & 3/5 & 尼洋河 + 古城氛围 & 可缩到 20 min \\
秀巴古堡 & 1.5/5 & 古碉楼历史 & \textbf{优先砍} \\
卡定沟 & 3.5/5 & 瀑布峡谷，和其他点有差异 & 45 min 快刷 \\
比日神山 & 2.5/5 & 山林/苯教氛围/观景 & 时间紧可砍 \\
色季拉 & 3--5/5 & 天气决定一切，晴天南迦巴瓦加分 & 云厚就 5 min 走 \\
扎西岗+贡措 & 4.5/5 & 近雪山、林海、田园 & 晴天尽量保 \\
易贡--尼屋穿越 & 5/5 景观 / 5/5 难度 & 连续地貌变化和峡谷 & 车辆/雨天不合适必须砍 \\
嘉乃玉措 & 4.5/5 & 开阔藏北湖景 & 保留 \\
麦地卡湿地 & 4/5 & 湿地、河流、藏北开阔感 & 保主体，东德错可砍 \\
东德错 & 3.5/5 & 湿地深处大湖、人少 & 时间/路况差先砍 \\
圣象天门 & 5/5 & 纳木错最佳近湖大尺度之一 & \textbf{整天核心，不砍} \\
蓬措 & 3/5 & 小众蓝湖 & 先砍 \\
那根拉 & 2/5 & 高海拔山口打卡 & 先砍 \\
羊八井野温泉 & 3.5/5 & 地热野景，差异化强 & 顺路时保留 \\
罗布林卡 & 2.5--4/5 & 宫苑/宗教人文，兴趣驱动 & 不懂人文者可砍 \\
西藏博物馆 & 5/5 & 系统理解西藏历史人文 & \textbf{拉萨人文首选} \\
羊湖 1--5 号 & 4.5/5 & 从高位到湖边连续变化 & 不要只留 1 号 \\
卡若拉 & 5/5 & 冰川近景与栈道爬升 & 保留 \\
恰央错 & 4.5/5 & 雪山脚下湖与空旷山谷 & 保留 \\
岗布沟 & 2.5--4.5/5 & 康布措+雪山；受栈道状况影响极大 & 先查开放，坏了就砍 \\
普莫雍措 & 4.5/5 & 高原蓝湖、藏南雪山远景 & 保留 \\
推瓦寺 & 2.5/5 & 人文/湖边小众点 & 前面超时就砍 \\
蒙达拉/库拉岗日 & 5/5 & 藏南雪山近景高潮 & 保留 \\
日托寺 & 3/5 & 小众寺庙与水上视角 & 赶机场时砍 \\
\bottomrule
\end{longtable}

% ==========================================================
\section{天气、季节与驾驶安全}

\subsection{藏东南的天气不能按“拉萨今天晴”推断}
巴松措、林芝、鲁朗、易贡、尼屋受暖湿气流影响更明显，夏季降水多，常在下午、傍晚或夜间出现阵雨/强对流。本次巴松措夜间遇到冰雹并直接影响第二天节奏。

执行规则：
\begin{itemize}
  \item 早晨优先完成徒步、雪山和非铺装支线；下午留主路和城市补给。
  \item 雨后不要立即冲新措、八盖--尼屋、东德错等支线。
  \item 看雪山（南迦巴瓦、加拉白垒、库拉岗日）不要只看“天气晴/雨”，更要看低云和山口云墙。
\end{itemize}

\subsection{高原驾驶的几个坑}
\begin{itemize}
  \item 长下坡：鲁朗--通麦尤其典型，控制速度、避免持续重踩刹车。
  \item 夜路：野生动物、牲畜、落石、无照明坑洼会把风险放大；本线偏远烂路尽量天黑前结束。
  \item 大车会车：林芝--鲁朗 G318 很窄，宁可停稳让车，不抢缝。
  \item 托底：过大坑不要“骑坑中间”，先观察坑深和车辙；必要时让副驾下车指挥。
  \item 胎壁：比“公路胎抓地”更重要的是抗冲击能力，备胎/补胎工具/充气泵必须确认可用。
  \item 燃油：偏远段看到可靠加油站且油量低于一半，通常就值得补，不要追求把油箱算到极限。
\end{itemize}

% ==========================================================
\section{实测费用：三人一次完整自驾的量级参考}

\mainbox{别拿这个表做报价单}{以下是 2026 年 8 月三人实际支出复盘，包含上海出发航班、成都中转、重庆短停等个人行程成本。机票、住宿、景区票价都会变化；这张表的意义是给\textbf{费用量级和结构}，不是保证未来同价。}

\begin{longtable}{p{4.8cm}r p{7.8cm}}
\toprule
项目 & 实付/记录 & 说明 \\
\midrule
机票 & 10775 元 & 三人，含上海--成都--拉萨、拉萨--重庆--上海 \\
租车 & 2600 元 & 长安深蓝 G318 \\
电费 & 576.1 元 & 多次快充/补电合计 \\
油费 & 528 元 & 增程兜底燃油 \\
餐饮 & 约 1895 元 & 含路餐和重庆下午茶；三人 \\
酒店 & 2967 元 & 巴松措、鲁朗/尼屋/那曲/当雄/羊湖/机场附近等 \\
景点 & 842 元 & 实际购买的景点票 \\
打车 & 149 元 & 主要城市/机场交通 \\
机场大巴 & 75 元 & 三人 \\
共享骑行 & 49 元 & 拉萨等城市短途 \\
补给 & 622.5 元 & 水、路餐、日用品等 \\
药品 & 215 元 & 本次个人准备，非通用处方 \\
停车 & 6.7 元 & 多数自然景点停车成本低 \\
\midrule
\textbf{总计} & \textbf{约 21300.3 元} & 三人全旅行总口径 \\
\textbf{扣除机票后} & \textbf{约 10525.3 元} & 更接近西藏地面部分 + 城市接驳的量级 \\
\bottomrule
\end{longtable}

\summarybox{车辆能源成本}{本次油+电合计约 \textbf{1104.1 元}。增程车的优势是主城区/县城用电降低成本，偏远段靠油确保不被充电桩绑架。}

% ==========================================================
\section{出门前清单：完整复制路线时至少带这些}

\subsection{证件与文件}
\begin{multicols}{2}
\begin{itemize}[label=\ck]
  \item 身份证原件
  \item 驾驶证
  \item 边境管理区通行证/电子证离线截图
  \item 学生证/优惠证件
  \item 租车合同、保险、救援电话
  \item 酒店与景区订单截图
  \item GPX 离线备份
  \item 高德 + OsmAnd 离线地图
\end{itemize}
\end{multicols}

\subsection{高原与个人健康}
\begin{multicols}{2}
\begin{itemize}[label=\ck]
  \item 血氧仪
  \item 个人常用药/长期药物 + 余量
  \item 医生建议的高原相关药物
  \item 便携氧气（到拉萨后购买）
  \item 电解质饮料/冲剂
  \item 止痛退烧、肠胃、晕车类常用药
  \item 润唇膏
  \item 保湿霜
  \item SPF50+ 防晒霜
  \item UV400 墨镜
\end{itemize}
\end{multicols}

\subsection{车辆与烂路用品}
\begin{multicols}{2}
\begin{itemize}[label=\ck]
  \item 备胎或确认原车补胎方案
  \item 充气泵
  \item 补胎工具
  \item 拖车钩/拖车绳（会正确使用再带）
  \item 强光手电/头灯
  \item 车载充电线
  \item 手机离线导航支架
  \item 垃圾袋
  \item 雨披
  \item 纸巾/湿巾
  \item 2--3 L 以上备用饮水
  \item 至少一顿完整路餐
\end{itemize}
\end{multicols}

\subsection{衣物}
\begin{multicols}{2}
\begin{itemize}[label=\ck]
  \item 防风外套/冲锋衣
  \item 抓绒或薄羽绒
  \item 防雨层
  \item 长裤
  \item 防滑徒步鞋
  \item 备用袜子
  \item 遮阳帽
  \item 保暖帽
  \item 薄手套
\end{itemize}
\end{multicols}

% ==========================================================
\section{临场决策树：遇到超时、下雨、高反，怎么砍最不亏}

\subsection{前两天超时}
\begin{enumerate}
  \item D1：先砍太昭长游，只短停；不加思金拉错；确保天黑前到巴松措。
  \item D2：若新措完整走，就不要再执着秀巴/比日神山；卡定沟根据出景区时间决定。
  \item 若林芝已晚：比日神山直接砍，林芝吃饭补给后翻色季拉；山口云厚不等南迦巴瓦。
\end{enumerate}

\subsection{易贡--尼屋前下雨}
\begin{enumerate}
  \item 先问通麦/易贡当地司机、酒店、施工人员路况。
  \item 若持续强降雨、塌方风险或夜间才进入八盖段：取消穿越。
  \item 不要因为“酒店已经订了尼屋”反向绑架安全决策。
\end{enumerate}

\subsection{那曲高反明显}
\begin{enumerate}
  \item 不去散步、夜游、跑步；保持休息和观察。
  \item 次日如果症状加重，不要继续冲圣象天门更高、更偏远区域。
  \item 出现严重危险表现，优先下降海拔和医疗，不把“完成路线”当目标。
\end{enumerate}

\subsection{最后一天卡若拉三景超时}
\begin{enumerate}
  \item 推瓦寺先砍。
  \item 日托寺直接砍。
  \item 保普莫 + 蒙达拉/库拉岗日；如果仍晚，洛扎住宿，第二天回贡嘎。
\end{enumerate}

% ==========================================================
\section{资料来源与时效说明}

\subsection{实测资料}
本文核心判断来自 2026 年 8 月一次完整自驾后的两份复盘：
\begin{itemize}
  \item 实际逐日行程、实际餐饮/酒店/景区/油电费用记录。
  \item 32 条补充注意事项，覆盖新措、易贡--尼屋、麦地卡、圣象天门、羊八井、羊湖、卡若拉、普莫--洛扎等路况和体验。
  \item OsmAnd 文件 \texttt{西藏大环线最终版.gpx}：1 条 route、1 条连续 track，67 个路线控制点、28628 个几何点；路线几何约 2608 km。该文件为\textbf{导航规划轨迹，不是 GPS 实际记录轨迹}。
\end{itemize}

\subsection{2026 年政策/官方信息核验}
\begin{itemize}
  \item 圣象天门自驾政策：西藏自治区文化和旅游厅转载班戈县公告，2026-06-15 起取消三日预约和每日 50 辆限制；官方购票核验后按指定线路自驾。\\
  \url{https://wlt.xizang.gov.cn/xccx/lytg/202606/t20260617_546017.html}
  \item 电子边境管理区通行证：2026-04-15 起可通过“移民局12367”线上申请。\\
  \url{https://wlt.xizang.gov.cn/xccx/lytg/202604/t20260415_535038.html}
  \item 西藏自治区边境管理条例：非本边境管理区户籍中国公民进入边境管理区，应持居民身份证和边境管理区通行证。\\
  \url{https://wjw.xizang.gov.cn/ztzl/wsfzxczl/202508/t20250808_494220.html}
  \item 西藏博物馆：个人可通过官方微信公众号预约，门票免费，个人可提前 7 天。\\
  \url{https://wlt.xizang.gov.cn/xccx/lytg/202603/t20260327_531698.html}
  \item Wilderness Medical Society 2024 高原病临床实践指南。\\
  \url{https://doi.org/10.1016/j.wem.2023.05.013}
\end{itemize}

\riskbox{时效声明}{\textbf{路况是 2026 年 8 月实测，不是永久事实。}西藏山区受降雨、塌方、施工、冬季积雪和临时管制影响极大；景区票价、开放、道路和边境管理要求也会变化。真正出发前，请重新确认天气、道路管制、景区公告和边境通行要求。}

% ==========================================================
\section{最后的路线结论}

如果只用一句话评价这条线：\textbf{风景非常强，多样性远高于“只去几个热门湖”的常规西藏路线，但它对海拔适应、车辆和临场决策的要求也明显更高。}

最值得多给时间的是：\textbf{巴松措/新措、鲁朗、易贡--尼屋穿越、嘉乃玉措/麦地卡、圣象天门、西藏博物馆、羊湖连续观景台、卡若拉--恰央错、普莫--蒙达拉--库拉岗日。}

最应该学会砍的是：\textbf{秀巴古堡、比日神山、那根拉、推瓦寺、日托寺，以及所有在天气/路况已经不合适时还想硬冲的支线。}

\summarybox{真正的核心原则}{\textbf{在西藏，“路能通”不等于“这段路好开”；“地图上只有 50 km”也不等于“一小时能到”。}把路况、海拔和天气都当作行程的一部分，而不是景点之间的空白，整条环线才会从“赶路”变成真正好玩的自驾。}

\end{document}
